\section{Деревья и их основные свойства}
\label{sec:trees-properties}

\subsection{Определение дерева}

Неориентированный граф называется лесом, если он не содержит циклов. Связный
лес называется деревом. Таким образом, дерево --- связный ациклический граф.
Компоненты связности леса являются деревьями.

Для графа $G=(V,E)$ с $n=|V|$ вершинами следующие утверждения эквивалентны:
\begin{enumerate}
    \item $G$ является деревом;
    \item $G$ связен и имеет $n-1$ ребро;
    \item $G$ ацикличен и имеет $n-1$ ребро;
    \item между любыми двумя вершинами существует единственная простая цепь;
    \item $G$ связен, но удаление любого ребра нарушает связность;
    \item $G$ ацикличен, но добавление любого нового ребра создаёт ровно один цикл.
\end{enumerate}

Из этих свойств следует, что каждое ребро дерева является мостом.

\subsection{Листья и степени вершин}

Вершина степени один называется листом, или висячей вершиной. Всякое дерево с
$n\geq2$ содержит не менее двух листьев. Это можно доказать, рассмотрев концы
максимальной простой цепи: у них не может быть дополнительных соседей, иначе
цепь можно было бы продолжить или возник бы цикл.

По лемме о рукопожатиях для дерева
\[
    \sum_{v\in V}d(v)=2(n-1).
\]
Если $n_k$ --- число вершин степени $k$, то
\[
    \sum_{k\geq1}n_k=n,
    \qquad
    \sum_{k\geq1}k n_k=2(n-1).
\]
Отсюда, например,
\[
    n_1=2+\sum_{k\geq3}(k-2)n_k.
\]

\subsection{Корневые деревья}

Если в дереве выделена вершина $r$, она называется корнем, а дерево ---
корневым. Для каждой вершины $v\neq r$ единственная цепь из $r$ в $v$
определяет родителя $p(v)$. Соседи, удалённые от корня на один уровень дальше,
называются детьми.

Уровень вершины --- расстояние от корня. Максимальный уровень называется высотой
дерева. Поддерево вершины $v$ состоит из $v$ и всех её потомков.

Ориентированное корневое дерево, или арборесценция, можно получить, направив все
рёбра от корня. Тогда полустепень захода удовлетворяет
\[
    d_{\mathrm{in}}(r)=0,
    \qquad
    d_{\mathrm{in}}(v)=1\quad(v\neq r),
\]
и каждая вершина достижима из корня.

\subsection{Упорядоченные и бинарные деревья}

Корневое дерево называется упорядоченным, если для детей каждой вершины задан
порядок. Бинарное дерево --- корневое дерево, в котором у каждой вершины не более
двух детей, различаемых как левый и правый.

Для полного бинарного дерева, где каждая внутренняя вершина имеет ровно двух
детей, число листьев $L$ и число внутренних вершин $I$ связаны равенством
\[
    L=I+1.
\]
Если высота бинарного дерева равна $h$ и корень имеет уровень $0$, то число
вершин не превосходит
\[
    1+2+\cdots+2^h=2^{h+1}-1.
\]

\subsection{Центр дерева}

Эксцентриситет вершины $v$ определяется как
\[
    e(v)=\max_{u\in V}d(u,v).
\]
Вершины с минимальным эксцентриситетом образуют центр дерева. Центр любого дерева
состоит либо из одной вершины, либо из двух смежных вершин. Его можно найти,
последовательно удаляя все листья: после нескольких слоёв останется центр.

Диаметр дерева --- длина максимальной простой цепи. Практический алгоритм его
нахождения:
\begin{enumerate}
    \item выполнить BFS или DFS из произвольной вершины и найти наиболее далёкую
          вершину $u$;
    \item выполнить обход из $u$; наиболее далёкая вершина $v$ является другим
          концом диаметра, а $d(u,v)$ --- длина диаметра.
\end{enumerate}

\subsection{Остовные деревья}

Остовным деревом связного графа $G$ называется подграф, содержащий все вершины
$G$ и являющийся деревом. Любой связный граф имеет остовное дерево: достаточно
выполнить DFS или BFS и взять рёбра, по которым впервые достигаются вершины.

Если рёбра имеют веса, задача минимального остовного дерева решается алгоритмами
Краскала или Прима. Для полного помеченного графа на $n$ вершинах число различных
остовных деревьев равно $n^{n-2}$ --- формула Кэли.

\subsection{Код Прюфера}

Помеченному дереву на вершинах $\{1,\ldots,n\}$ соответствует последовательность
Прюфера длины $n-2$. На каждом шаге удаляют лист с наименьшей меткой и записывают
метку его соседа. Код задаёт дерево однозначно и даёт простое доказательство
формулы Кэли, поскольку существует ровно $n^{n-2}$ последовательностей длины
$n-2$ над алфавитом из $n$ меток.

\subsection{Что важно сказать на экзамене}

Главная характеристика дерева --- эквивалентность связности с отсутствием циклов
при $n-1$ рёбрах. Нужно уметь объяснить единственность пути, наличие листьев,
понятия корня, уровня и поддерева, а также связь деревьев с остовами графа.

\newpage